\documentclass[11pt]{article}
\usepackage{amsmath,amssymb,booktabs,geometry,xcolor}
\geometry{margin=1.1in}
\newcommand{\Tr}{\mathrm{Tr}}
\newcommand{\Sig}{\Sigma}
\title{DTXQCD notes: dynamical free-field Fierz verification\\
and the redundancy-resolved (traceless) auxiliary action}
\author{}
\date{2026-07-03}

\begin{document}
\maketitle

\section{Exact integration of the fermion-blind singlet mode}
\label{sec:traceless}

\subsection{Setup and the $u$-universality input}

The v2 DTXQCD auxiliary fields comprise the color--flavor $6\times 6$
Hermitian matrices $\sigma,\pi$ (and the complex-symmetric $d,n$, which play
no role here) plus the real singlets $s,p$, with Gaussian prior
\begin{equation}
S_{\rm aux} \;=\; \frac{\lambda^{2}}{2}\sum_x\Big[\Tr\sigma^{2}+\Tr\pi^{2}
 +\Tr d^{2}+\Tr n^{2}+s^{2}+p^{2}\Big].
\label{eq:legacyprior}
\end{equation}
The doubled fermion matrix $M_{48}$ depends on $(\pi,p)$ only through the
combination $\pi + p\,\mathbb{1}_{6}$ (and on $(\sigma,s)$ through
$\sigma + s\,\mathbb{1}_{6}$).  Writing
$\pi=\tilde\pi+\pi_{0}\mathbb{1}_{6}$ with $\pi_{0}=\Tr\pi/6$,
$\Tr\tilde\pi=0$, the matrix depends on the singlet sector only through
\begin{equation}
u \;=\; \pi_{0}+p ,
\end{equation}
with \emph{unit relative weight}.  This is not an assumption: the code gate
verifies $M(\pi+c\,\mathbb{1},\,p)=M(\pi,\,p+c)$ to $6\times10^{-17}$
(host operator), $S$-level equality to $10^{-16}$ relative for the LogDet
and rational actions, and the fermionic force identity
$\partial S_{f}/\partial(\Tr\sigma) \big/ \partial S_{f}/\partial s
 = 1.000000$ for both LogDetEO and RationalFull.  The orthogonal combination
$w\propto \pi_{0}-p/6$ is exactly fermion-blind: a free Gaussian mode with
prior mass only, responsible for the slow $\tau_{\rm int}\simeq 2$
trajectories observed in $\langle p\rangle$ at $\lambda=6$ production and
for frozen-mode pathologies at small $\lambda$.

\subsection{Marginalization and the $6/7$ coefficient}

Under the decomposition, $\Tr\pi^{2}=\Tr\tilde\pi^{2}+6\pi_{0}^{2}$, so the
singlet-sector prior is
\begin{equation}
S \;\supset\; \tfrac12\,a\,\pi_{0}^{2}+\tfrac12\,b\,p^{2},
\qquad a=6\lambda^{2},\quad b=\lambda^{2}.
\end{equation}
The trace mode is six times stiffer than $p$ because it appears six times in
$\Tr\pi^{2}$.  Since the physics couples only to $u=\pi_{0}+p$, the blind
direction may be integrated out exactly; the variances add,
\begin{equation}
\mathrm{Var}(u)=\frac1a+\frac1b=\frac{7}{6\lambda^{2}}
\;\;\Longrightarrow\;\;
S_{\rm eff}(u)=\frac{u^{2}}{2\,\mathrm{Var}(u)}=\frac{3}{7}\lambda^{2}u^{2}.
\label{eq:seff}
\end{equation}
The exactly equivalent single-field theory therefore has $\tilde\pi$
(traceless) with unchanged coefficient $\lambda^{2}/2$, and $p$ carrying the
full $u$-mode with
\begin{equation}
\boxed{\;S_{p}=\frac{3}{7}\lambda^{2}p^{2}
 =\frac{\lambda^{2}}{2}\cdot\frac{6}{7}\,p^{2}\;}
\end{equation}
i.e.\ a coefficient rescale of $6/7$ (identically for $\sigma/s$).  Every
fermionic observable is preserved exactly: the joint law of
$(\tilde\pi,u)$ --- everything $M_{48}$ sees --- is unchanged.  Because the
transformation is a linear rotation of a Gaussian followed by an exact
marginalization, there is no Jacobian and the Metropolis test remains exact.

Consistency checks:
\begin{itemize}
\item \emph{Saddle map.}  With a source $J$ coupling to $u$: legacy
$\pi_{0}=J/6\lambda^{2}$, $p=J/\lambda^{2}$, $u=7J/6\lambda^{2}$; new
$u=J/[(6/7)\lambda^{2}]=7J/6\lambda^{2}$.  Identical.  Equivalently
$\langle s\rangle_{\rm new}=\tfrac{7}{6}N_{f}\Sig/\lambda^{2}
=\langle \Tr\sigma/6+s\rangle_{\rm legacy}$, and
$\langle\Tr\sigma\rangle\equiv 0$ by construction.
\item \emph{Why the 2026-06-15 attempt failed.}  That version used
$\lambda^{2}/4$, i.e.\ $\mathrm{Var}(u)=2/\lambda^{2}$ instead of
$7/6\lambda^{2}$: the $u$ fluctuations were $71\%$ too wide in variance.
Its \emph{saddle} was self-consistently absorbed ($\langle s\rangle$
doubled, which the scout logged as success) --- Gaussian means are
forgiving; widths are not --- and the Fierz mass scan falsified it.
\item \emph{Production cross-check.}  The interacting $16^{3}\times48$
ensembles at $\lambda=5,6$ satisfy the wall-normalized identity
$\langle\Tr\sigma\rangle=\langle s\rangle$ to $0.04\%/0.16\%$ and the
saddle scaling $\langle s\rangle\propto 1/\lambda^{2}$
($1007/701\simeq(6/5)^{2}$), confirming the assumed mode structure in situ.
\end{itemize}

\subsection{Implementation}

Enabled by \texttt{DTXQCD\_TRACELESS\_AUX=1} (default off = legacy):
$\sigma,\pi$ are projected traceless at momentum generation, force
projection, and \texttt{Project()} via
$X\to X-\tfrac16(\Tr X)\mathbb{1}$ (device-resident expression template;
unit-tested exact, idempotent, off-diagonal preserving); the $s,p$ action
and force carry the factor $6/7$; the initial draw uses width
$\sqrt{7/6}/\lambda$ for $s,p$ and drops the diagonal $\sigma$ saddle shift
(it \emph{is} the trace mode), placing
$\langle s\rangle=\tfrac76 N_{f}\Sig/\lambda^{2}$.

\subsection{Dynamical equivalence validation}

Free-field HMC (Sec.~\ref{sec:fierz} settings, auto-sized rational
bracket), knob ON vs OFF at $\lambda=10$:
\begin{center}
\begin{tabular}{lcc}
\toprule
$\Sig_{\rm DTX}/\Sig_{W}$ & traceless-aux (ON) & legacy (OFF)\\
\midrule
$m=0.1$ & $1.00126(81)$ & $1.00207(83)$\\
$m=1$   & $1.00120(75)$ & $1.00157(72)$\\
\bottomrule
\end{tabular}
\end{center}
The two actions agree with each other ($\le0.7\sigma$ per mass) and each
run is consistent with unity by the test's criteria;
$\langle\Tr\sigma\rangle=\mathcal{O}(10^{-19})$ dynamically under the knob,
and all analytic-vs-finite-difference force ratios $=1.000000$ per action
and per aux slot.  (Pooling the four corrected-bracket runs leaves a
common $+0.15(4)\%$ offset --- $\sim\!10\times$ below the wrong-bracket
signal and action-independent; plausible residuals at this level include
the $10^{-8}$ rational tolerance and numerator--denominator noise
correlations in the ratio estimator.  A high-statistics run with a
deliberately over-wide bracket would pin it down if sub-per-mille
exactness verification is ever needed.)

The production streams provide the same verification in situ: the
$\lambda=6,5,4$ traceless-aux chains reproduce the legacy saddle map
$\langle s\rangle_{\rm TLA}=\tfrac76\langle s\rangle_{\rm legacy}
=\langle \Tr\sigma/6+s\rangle_{\rm legacy}$ to $0.02$--$0.04\%$, hold
$\langle\Tr\sigma\rangle=\langle\Tr\pi\rangle\equiv0$ to $10^{-16}$
through live HMC, and collapse the blind-mode autocorrelation
($\tau_{\rm int}(p)$: $1.8\to0.4$ at $\lambda=6$; legacy $\lambda=4$ had
$\tau_{\rm int}(p)=6.7$, the fleet's worst).  (The wrong-bracket runs of
Sec.~\ref{sec:fierz} also satisfied ON$\,=\,$OFF at every mass --- both
actions sample the same distribution whatever the rational weight is ---
so the equivalence conclusion never depended on the bracket.)

\section{Dynamical free-field verification of Fierz exactness}
\label{sec:fierz}

\subsection{The statement and the estimator}

The DTXQCD construction must satisfy
$\langle\mathcal{O}\rangle_{\rm DTXQCD}=\langle\mathcal{O}\rangle_{\rm QCD}$
\emph{exactly} at every $\lambda$, $V$, $m$, $\beta$: the four-quark cross
terms cancel by an exact Fierz identity, and only a failure to sample the
correct distribution can break the equality.  The free-field test fixes
$U=\mathbb{1}$ (frozen gauge), runs dynamical-auxiliary HMC with the full
rational pseudofermion weight $|\det M_{48}|^{1/2}$
(\texttt{USE\_FULL\_PF=1}), and measures the per-quark stochastic trace
$\Sig_{\rm DTX}=\langle\Tr M_{48}^{-1}\rangle/(2N_{f})$ against the exact
free Wilson--clover reference $\Sig_{W}$ on the same geometry ($4\times$
noise statistics on the reference so its error is subdominant).  The
per-measurement ratio series is logged, enabling Madras--Sokal
$\tau_{\rm int}$-corrected errors; at measurement stride 5 the series is
statistically uncorrelated ($\tau_{\rm int}$ consistent with $0.5$).

\subsection{Settings}

\begin{center}
\begin{tabular}{ll}
\toprule
lattice & $4^{3}\times 8$ (volume check: $8^{3}\times16$, $8^{3}\times32$)\\
gauge & $U=\mathbb{1}$, frozen; $c_{\rm sw}=1.24930970916466$\\
masses & $m\in\{100,\,10,\,1,\,0.1\}$ (heavy $\to$ light scan)\\
couplings & $\lambda\in\{20,\,10,\,5\}$\\
action & traceless-aux (Sec.~\ref{sec:traceless}); legacy for ON/OFF pairs\\
HMC & MN2, trajL $=1$, \textbf{MDSTEPS $\ge 1.5\,\lambda\sqrt2\,$trajL}\\
 & (16 at $\lambda\!=\!10$, 8 at $\lambda\!=\!5$, 32 at $\lambda\!=\!20$)\\
statistics & $N_{\rm therm}=40$, $N_{\rm prod}=60$, stride 5 ($N=13$ samples)\\
rational & $x^{-1/4}$, bracket \textbf{auto-sized from $m$}
 (see caveat below), tol $10^{-8}$\\
\bottomrule
\end{tabular}
\end{center}

Two harness pitfalls, both load-bearing:
\begin{enumerate}
\item \emph{Symplectic stability.}  The auxiliary prior gives the singlet
modes frequency $\omega=\lambda\sqrt2$; leapfrog requires
$\omega\varepsilon<2$.  The historical default (MDSTEPS $=4$, trajL $=1$)
sits at $\omega\varepsilon=3.5$ for $\lambda=10$: $dH\simeq+240$ every
trajectory, $100\%$ rejection, and the ensemble silently frozen at its
initialization --- all pre-2026-07-03 runs of this gate sampled statics,
and their ``saddle passes'' were init tautologies.
\item \emph{The rational bracket (the trap).}  The heavy-mass variant of
the test hardcoded its Remez window to $[10^{4},10^{7}]$ (sized for
$m=1000$).  At light mass the $M_{48}^{\dagger}M_{48}$ spectrum sits up to
six decades \emph{below} the window.  A rational approximation used
self-consistently outside its bracket samples a wrong---but
Metropolis-consistent---distribution: an action-level error that no
accept/reject step can filter.  The resulting spurious signal mimicked
physics disturbingly well (next subsection).
\end{enumerate}

\subsection{The spurious signal and its execution}

With the wrong bracket, the dynamical ratio showed a ``violation'' with
seductive systematics --- growing toward light mass, scaling
$\sim1/\lambda^{2}$, volume-stable:
\begin{center}
\begin{tabular}{lcccc}
\toprule
 & $m=100$ & $m=10$ & $m=1$ & $m=0.1$\\
\midrule
$\lambda=20$ & --- & --- & --- & $+0.18(8)\%$\\
$\lambda=10$ & $+0.13(7)\%$ & $+0.21(7)\%$ & $+1.15(8)\%$ & $+1.70(9)\%$\\
$\lambda=5$  & --- & --- & --- & $+7.50(11)\%$\\
\bottomrule
\end{tabular}
\end{center}
($m=100$ --- the one point \emph{inside} the bracket --- was the only clean
one; volume scan at $\lambda=5$, $m=0.1$: $+7.51(11)\%$, $+7.70(3)\%$,
$+7.74(2)\%$ at $4^{3}\times8$, $8^{3}\times16$, $8^{3}\times32$.)

Three independent tests dispatched it:
\begin{enumerate}
\item \emph{Sign problem} (the a priori favorite, since only sampling can
break the identity): dense $\gamma_{5}M_{48}$ spectra over 12 prior draws
at each $\lambda\in\{20,10,5,4\}$ at the test point gave
$\mathrm{sign}\,\mathrm{Pf}=+1$ in $48/48$ cases with spectral gap
$\ge0.28$ --- the free-field $m=+0.1$ point is sign-clean (positive Wilson
mass, no real-mode crossings).  Ruled out.
\item \emph{Finite-volume zero mode}: the excess is volume-flat over a
$32\times$ range while sign damage would compound exponentially and
zero-mode effects would shrink as $1/\sqrt V$.  Ruled out (both ways).
\item \emph{The kill shot}: same point ($\lambda=5$, $m=0.1$,
$4^{3}\times8$) with bracket $[10^{-4},10^{2}]$, degree 24:
\begin{equation}
\boxed{\;\Sig_{\rm DTX}/\Sig_{W} \;=\; 0.99946(109)\;}
\end{equation}
The $66\sigma$ ``violation'' collapses to unity at $0.5\sigma$.
\end{enumerate}

\subsection{Conclusions}

Fierz exactness holds \emph{dynamically} in the free-field theory at the
most sensitive point probed (sub-per-mille at $\lambda=5$, $m=0.1$), for
both the legacy and traceless-aux actions.  The production ensembles were
never exposed: the generator's \texttt{RAT\_AUTO\_HI=1} rebuilds the
rational bracket from a Lanczos estimate of the operator spectrum at every
refresh (Sec.~\ref{sec:remez}).  The operational lesson: \textbf{when a
DTXQCD observable deviates, check the Remez window against the actual
spectrum first} --- it is the cheapest sampling failure to create and the
cheapest to detect.

\section{Remez auto-tuning}
\label{sec:remez}

The DTXQCD rational actions hold two Zolotarev/Remez approximations on a
window $[\ell,h]$: $x^{-1/4}$ (used by $S$ and the force multishift ---
one square root for the $N_f=2$ degeneracy, one for
$\det M_{48}\to\mathrm{Pf}$) and $x^{+1/8}$ (used by \texttt{refresh}).
The generators historically hardcoded $[\ell,h]=[0.1,64]$.  Two coupled
failure modes appear whenever the spectrum of the inverted operator leaves
the window:
\begin{itemize}
\item \emph{Below $\ell$:} the approximation over-extrapolates.  The
per-shift multishift tolerances become effectively unattainably tight (CG
stalls), the pseudofermion force acquires large spurious contributions
($dH$ blowup) --- and, more insidiously, if the run \emph{does} proceed,
the sampled weight is a smooth wrong function of the low modes: the
distribution is mis-sampled with no Metropolis signal whatsoever
(Sec.~\ref{sec:fierz}).
\item \emph{Above $h$:} the hard end of the spectrum is misweighted and
the Remez error bound no longer applies.
\end{itemize}

\texttt{DTXQCDRemezAutoScale} (env master switch \texttt{RAT\_AUTO\_HI=1};
byte-identical to the fixed window when off) removes the exposure: at each
\texttt{refresh()} a short plain Lanczos (\texttt{RAT\_LANCZOS\_NM}$=30$
Krylov vectors; Eigen tridiagonal solve only) estimates
$\lambda_{\min},\lambda_{\max}$ of \emph{the operator the multishift
actually inverts} --- $M_{\rm pc}^{\dagger}M_{\rm pc}$ on the odd
checkerboard for the EO action, $M^{\dagger}M$ on the full grid for the
non-EO action --- and the Remez pair is rebuilt whenever the safety-scaled
window $[\,s_{\ell}\lambda_{\min},\,s_{h}\lambda_{\max}\,]$
($s_{\ell}=0.5$, $s_{h}=1.5$ by default) escapes the current one.  Matching
the Lanczos operator to the inverted operator is essential: the bounds must
be in the same units as the shifts.  Production has run with the
auto-scaler on since 2026-07-02, which is why the ensembles were never
exposed to the Sec.~\ref{sec:fierz} pathology; the free-field harness,
which constructed its rational once in the constructor from a fixed
(m=1000-sized) window, was.  Post-mortem policy: any fixed-bracket harness
must either auto-size the window from the run parameters (the free-field
test now derives $[\ell,h]$ from the input mass) or assert
$s_{\ell}\lambda_{\min}\ge\ell$ at run time.

\section{Auxiliary-field initialization (\texttt{AUX\_INIT})}
\label{sec:auxinit}

Fresh starts draw the auxiliary fields from their Gaussian prior (width
$1/\lambda$ per site; $\sqrt{7/6}/\lambda$ for $s,p$ under
\textsc{traceless\_aux}) around a self-consistent saddle.  With
\texttt{AUX\_INIT\_AUTO=1} the saddle condensate $\Sig^{*}$ solves
$\Sig=\Sig_{\rm DTXQCD}(\Sig)$, where $\Sig_{\rm DTXQCD}$ is the per-quark
stochastic trace $\Tr M_{48}^{-1}/(2N_{f})$ evaluated on a \emph{mean-only}
fill (fluctuations zeroed: the saddle is a mean-field quantity, and at
small $\lambda$ full-width fluctuations drive $M_{48}$ near-singular and
stall the trace CG).  The solver is Picard-first,
$\Sig_{n+1}=\Sig_{\rm DTXQCD}(\Sig_{n})$ from $\Sig_{0}=0$: the Jacobian
scales as $N_{f}/\lambda^{2}$, so at production couplings it contracts in
one or two evaluations (at $\lambda=10$ the saddle is mean-field to
$<0.1\%$); a bisection fallback covers the non-contracting small-$\lambda$
regime.  Tolerance \texttt{AUX\_INIT\_TOL} defaults to $0.1$ --- ample,
since the first trajectories equilibrate the residual.

\emph{Interaction with \textsc{traceless\_aux}.}  No explicit factors are
needed anywhere in the solver: the $u$-mode entering $M_{48}$ is identical
in the two conventions for the same $\Sig$ (that is the content of the
equivalence), so $\Sig_{\rm DTXQCD}(\Sig)$ --- and hence $\Sig^{*}$ --- is
convention-independent.  All convention dependence is localized in the
final field placement inside \texttt{FillAuxFields}: legacy splits the
saddle as $\sigma_{0}=\Sig/(N_{c}\lambda^{2})$ diagonal shift plus
$s=N_{f}\Sig/\lambda^{2}$; traceless-aux drops the $\sigma$ shift (it
\emph{is} the trace mode) and places
$s=\tfrac{7}{6}N_{f}\Sig/\lambda^{2}$, reproducing the same $u$.  (A stale
log label ``$\langle s\rangle^{*}=2N_{f}\Sig/\lambda^{2}$'' --- a leftover
of the reverted 2026-06-15 convention that matched neither action --- was
corrected in passing.)  Explicit alternatives: \texttt{AUX\_INIT=}$\Sig$
skips the solve; \texttt{ZERO\_ALL\_AUX=1} gives the cold start used by
frozen-aux controls.

\section{Locating the sign-problem onset from random auxiliary draws}
\label{sec:signboundary}

\subsection{Method}

HMC samples $|\mathrm{Pf}(K M_{48})|\,e^{-S_{\rm aux}}$; physics requires
reweighting by $\mathrm{sign}\,\mathrm{Pf}$, which is feasible only where
$\langle\mathrm{sign}\rangle$ is not exponentially small.  A cheap,
HMC-free diagnostic of \emph{where} the sign turns on: fix a gauge
configuration, draw the auxiliary fields from their Gaussian prior (width
$1/\lambda$, the same marginal HMC equilibrates to up to
$|\mathrm{Pf}|$-weighting), and compute the signed Pfaffian densely for
each draw.  Concretely, $\gamma_{5}M_{48}$ is Hermitian; its full spectrum
is obtained by cuSOLVER \texttt{Zheevd} (dimension $48V$; $V=4^{3}\times8$
$\Rightarrow$ $24576$).  The charge-conjugation doubling forces exact
Kramers pairing of the eigenvalues (measured pair splitting
$\lesssim3\times10^{-14}$), so
\begin{equation}
\mathrm{sign}\,\mathrm{Pf}(KM_{48})=(-1)^{n_{-}/2},
\end{equation}
with $n_{-}$ the number of negative eigenvalues; the same diagonalization
yields $\log|\det|$ for $|\mathrm{Pf}|$-weighted averages and the spectral
gap $\min|\mathrm{eval}|$, which measures the distance to the nearest sign
flip.  Because the prior is \emph{wider} than the
$|\mathrm{Pf}|$-reweighted equilibrium (the weight suppresses
near-singular configurations), prior draws give a conservative
(pessimistic) bound on sign mixing: a clean prior implies clean dynamics.

\subsection{Results: the boundary at the production mass}

Tepid gauge ($U=\exp(0.01$-Lie noise$)$; the weak-field $0.1$ amplitude
gives statistically identical results), $m=-0.245$,
$c_{\rm sw}=1.2493$, $4^{3}\times8$, 12 draws per $\lambda$:

\begin{center}
\begin{tabular}{lccc}
\toprule
$\lambda$ & signs $(+/-)$ & median $\min|\mathrm{eval}|$ & verdict\\
\midrule
$0.1$--$2$ & $\sim$6/6 (coin flip) & $0.000$--$0.012$ & gapless, fully mixed\\
$3$   & 6/6            & $0.011$--$0.015$ & mixed\\
$4$   & \textbf{12/12 $+$} & $0.40$ & clean\\
$5$   & 12/12 $+$      & $0.52$ & clean\\
$6$   & 12/12 $+$      & $0.56$ & clean\\
\bottomrule
\end{tabular}
\end{center}

The onset is sharp: between $\lambda=3$ and $\lambda=4$ the spectral gap
jumps by a factor $\sim\!25$ and the sign goes from a fair coin to
unanimous.  $|\mathrm{Pf}|$-weighting of the same draws does not move the
boundary.  The mechanism is the auxiliary prior width $1/\lambda$ dragging
the (unprotected, Wilson) real modes of the near-critical operator through
zero; consistently, the boundary is parameter-dependent in exactly the
expected way:
\begin{itemize}
\item \emph{Gauge roughness}: tepid vs.\ weak-field identical; hot
(random) gauge is mixed at least through $\lambda=2$ and saturates the
diagnostic --- the onset should be probed on gauge fields representative
of the target ensemble.
\item \emph{Quark mass}: at $m=+0.1$, $U=\mathbb{1}$ (the free-field
Fierz point of Sec.~\ref{sec:fierz}), all $48$ draws over
$\lambda\in\{20,10,5,4\}$ are $+1$ with gap $\ge0.16$: a positive Wilson
mass has no real-mode crossings to drive.  The boundary is a property of
the near-critical mass region, not of the construction per se.
\item \emph{Discretization}: the same probe applied to TXQCD's naive
$s$-broadcast M\"obius DWF coupling flips \emph{earlier} (7/8 draws
mixed already at $\lambda=3$) --- the naive five-dimensional coupling
worsens the sign problem, so a DWF cure requires chirally sandwiched
(overlap-type) couplings.
\end{itemize}

\subsection{Volume scaling and the viable window}

Sign violations are spatially local (near-zero modes localize on
$\mathcal{O}(1)$ regions), so to good approximation patches of volume
$V_{0}$ flip independently and
\begin{equation}
\langle\mathrm{sign}\rangle_{V}\;\sim\;
\langle\mathrm{sign}\rangle_{V_{0}}^{\,V/V_{0}} .
\end{equation}
With the measured $\mathcal{O}(1)$ per-$4^{3}\times8$-patch mixing at
$\lambda\le3$, reweighting at $16^{3}\times48$ ($V/V_{0}=384$) is
exponentially dead; conversely the unanimous-plus region needs
$\langle\mathrm{sign}\rangle_{V_{0}}$ so close to 1 that even large
volumes stay clean --- the per-trajectory
$\gamma_{5}M_{48}$ low-eigenvalue monitor in production (the
\texttt{g5M} observer) certifies this continuously: no gap
collapse $\Rightarrow$ no sign flips.  These measurements fix the DTXQCD
production window to $\lambda\ge4$ at the $m=-0.245$ point, with
$\lambda=5$--$6$ preferred (larger gap, cheaper solves).  One structural
caveat for EO monitoring: the even-site-local Pfaffian factor and the Schur
factor flip \emph{together} at $M_{ee}$ crossings (a factorization
artifact, verified at $\lambda=4$), so the EE-block sign alone is not the
Pfaffian sign --- sign certification must use the full-operator spectrum.

\end{document}
